\documentclass{article}
\usepackage{amsmath, amssymb, amsthm}
\usepackage{hyperref}
\usepackage{geometry}
\geometry{margin=1in}

\title{Erd\H{o}s Problem \#642}
\date{Last updated: 2025-08-31}
\author{Source: erdosproblems.com}

\begin{document}
\maketitle

\noindent\textbf{Status:} OPEN \quad \textbf{No prize}

\noindent\textbf{Tags:} graph theory, cycles

\medskip
\noindent\textbf{Problem Statement:}

Let $f(n)$ be the maximal number of edges in a graph on $n$ vertices such that all cycles have more vertices than chords. Is it true that $f(n)\ll n$?

\bigskip
\noindent\textbf{Remarks:}

A problem of Hamburger and Szegedy. A chord is an edge between two vertices of the cycle which are not consecutive in the cycle.

Chen, Erd\H{o}s, and Staton \cite{CES96} proved $f(n) \ll n^{3/2}$. Dragani\'{c}, Methuku, Munh\'{a} Correia, and Sudakov \cite{DMMS24} have improved this to\[f(n) \ll n(\log n)^8.\]

\bigskip
\noindent\textbf{References:}

[CES96] Chen, Guantao and Erd\H{o}s, Paul and Staton, William, Proof of a conjecture of {B}ollob\'{a}s on nested cycles. J. Combin. Theory Ser. B (1996), 38--43.
 
 [DMMS24] Dragani\'c, Nemanja and Methuku, Abhishek and Munh\'{a}{}
Correia, David and Sudakov, Benny, Cycles with many chords. Random Structures Algorithms (2024), 3--16.

\bigskip
\noindent\small{Source: \url{https://www.erdosproblems.com/642}}
\end{document}
